\section*{Chapter \ref{ch:g10:surplus-and-efficiency} --- Surplus and Efficiency}
% figdata: figdata/grade-10/08-surplus-and-efficiency/out/digital_wta.csv (cited, ledger g10:bce2019-*); Brenland numbers invented.

\begin{solution}{exo:g10:surplus-and-efficiency:1}
$40 - 32 = 8$ dalens. At 45 dalens she does not buy, and her surplus is 0.
\end{solution}

\begin{solution}{exo:g10:surplus-and-efficiency:2}
$20 - 12 = 8$ dalens. He would not sell below 12 dalens.
\end{solution}

\begin{solution}{exo:g10:surplus-and-efficiency:3}
Search engines, \usd{17530}: about $17530/322 \approx 54$ times the median for social media (\usd{322}). Music (\usd{168}) and messaging (\usd{155}) were
closest.
\end{solution}

\begin{solution}{exo:g10:surplus-and-efficiency:4}
At $p = 10$, $Q = 40$; nobody buys above 30. \omterm{def:g10:surplus-and-efficiency:consumer-surplus}{Consumer surplus} $\tfrac12 \times 40 \times (30 - 10) = 400$ dalens.
\end{solution}

\begin{solution}{exo:g10:surplus-and-efficiency:5}
At $p = 15$, $Q = 30$; nothing is offered below 5. \omterm{def:g10:surplus-and-efficiency:producer-surplus}{Producer surplus} $\tfrac12 \times 30 \times (15 - 5) = 150$ dalens.
\end{solution}

\begin{solution}{exo:g10:surplus-and-efficiency:6}
$100 - 2p = 3p - 60$ gives $5p = 160$, $p^* = 32$, $Q^* = 36$. \omterm{def:g10:surplus-and-efficiency:consumer-surplus}{Consumer surplus} $\tfrac12 \times 36 \times (50 - 32) = 324$; \omterm{def:g10:surplus-and-efficiency:producer-surplus}{producer surplus} $\tfrac12
\times 36 \times (32 - 20) = 216$; total 540 (check: $\tfrac12 \times 36 \times (50 - 20) = 540$).
\end{solution}

\begin{solution}{exo:g10:surplus-and-efficiency:7}
$80 - p = p$ gives $p^* = 40$, $Q^* = 40$. \omterm{def:g10:surplus-and-efficiency:consumer-surplus}{Consumer surplus} $\tfrac12 \times 40 \times (80 - 40) = 800$; \omterm{def:g10:surplus-and-efficiency:producer-surplus}{producer surplus} $\tfrac12 \times 40 \times 40 = 800$;
total 1600. The two lines are equally steep, so the triangles above and below the price are mirror images of each other.
\end{solution}

\begin{solution}{exo:g10:surplus-and-efficiency:8}
At 18 units the demand price is $(100 - 18)/2 = 41$ and the supply price $(18 + 60)/3 = 26$: a gap of 15. The lost triangle is $\tfrac12 \times (36 - 18)
\times 15 = 135$, so \omterm{def:g10:surplus-and-efficiency:total-surplus}{total surplus} is $540 - 135 = 405$ dalens.
\end{solution}

\begin{solution}{exo:g10:surplus-and-efficiency:9}
$120 - 4p = 2p - 18$ gives $6p = 138$, $p^* = 23$, $Q^* = 28$. \omterm{def:g10:surplus-and-efficiency:consumer-surplus}{Consumer surplus} $\tfrac12 \times 28 \times (30 - 23) = 98$ (was 72); sellers now offer from 9
dalens, so \omterm{def:g10:surplus-and-efficiency:producer-surplus}{producer surplus} $\tfrac12 \times 28 \times (23 - 9) = 196$ (was 144). Both sides gain: buyers 26 and sellers 52 dalens.
\end{solution}

\begin{solution}{exo:g10:surplus-and-efficiency:10}
The staircase has the steps 10, 9, 8, 7, 6, 4, 3, 2, 1. At 5 dalens the five buyers with values 10, 9, 8, 7 and 6 buy, and the \omterm{def:g10:surplus-and-efficiency:consumer-surplus}{consumer surplus} is $5 + 4 + 3
+ 2 + 1 = 15$ dalens.
\end{solution}

\begin{solution}{exo:g10:surplus-and-efficiency:11}
$126 - 4p = 2p - 24$ gives $6p = 150$, $p^* = 25$, $Q^* = 26$. Nobody buys above 31.5: \omterm{def:g10:surplus-and-efficiency:consumer-surplus}{consumer surplus} $\tfrac12 \times 26 \times 6.5 = 84.5$ (was 72);
\omterm{def:g10:surplus-and-efficiency:producer-surplus}{producer surplus} $\tfrac12 \times 26 \times (25 - 12) = 169$ (was 144). Sellers gained more: 25 against 12.5.
\end{solution}

\begin{solution}{exo:g10:surplus-and-efficiency:12}
One would need to show that some buyers who do not buy value the good above what some sellers who do not sell would need to produce it: for example a persistent
queue or waiting list at a price fixed below equilibrium, or firms that would produce at the going price but are kept out. \omterm{def:g10:surplus-and-efficiency:wtp}{Willingness to pay} can be learned
from what people pay elsewhere or in a black market; costs from what similar firms charge where they are allowed to sell. A controlled experiment, as in a
laboratory market, measures both directly.
\end{solution}

\begin{solution}{exo:g10:surplus-and-efficiency:13}
\omterm{def:g10:surplus-and-efficiency:total-surplus}{Total surplus} with $Q^*$ units is the triangle between the two lines from 0 to $Q^*$. With $Q$ units it is the part of that triangle to the left of the vertical
line at $Q$. The rest is a triangle with a vertical side of length $g$ (from the supply price to the demand price at $Q$) and its third corner at the equilibrium
point, a horizontal distance $Q^* - Q$ away: its area is $\tfrac12 (Q^* - Q)\,g$, positive whenever $Q < Q^*$. For $Q > Q^*$ the extra units have a demand price
below the supply price, and the same reasoning subtracts a triangle of area $\tfrac12 (Q - Q^*)\,g$. Either way \omterm{def:g10:surplus-and-efficiency:total-surplus}{total surplus} is smaller than at $Q^*$.
\end{solution}

\begin{solution}{exo:g10:surplus-and-efficiency:14}
A rich person would pay 100 dalens for a concert ticket and a poor student 40, although the student would enjoy it more. Moving the ticket from the rich person
to the student lowers \omterm{def:g10:surplus-and-efficiency:total-surplus}{total surplus} by 60 dalens, yet the two together may well be better off. \omterm{def:g10:surplus-and-efficiency:total-surplus}{Total surplus} adds money values, which depend on income; it is a
precise measure of what people would pay, not of how much they gain in well-being, and judging a market by it gives more weight to people with more money.
\end{solution}

\begin{solution}{exo:g10:surplus-and-efficiency:15}
Each hundred eggs costs the neighbours $50 \times 2 = 100$ dalens, a cost in no seller's supply curve. The market trades every hundred eggs whose value to buyers
exceeds the farm's own cost, including some whose value is below the farm's cost plus the neighbours' 100 dalens: those units lower the surplus of everyone
concerned. The quantity that maximises it is smaller.
\end{solution}

\begin{solution}{pb:g10:surplus-and-efficiency:1}
\textbf{1.} Demand: steps at 30, 28, 25, 22, 20, 18, 15, 12, 10, 8. Supply: steps at 6, 8, 11, 13, 16, 19, 23, 26.

\textbf{2.} At 21: 4 wanted (30, 28, 25, 22), 6 offered (6 to 19). At 17: 6 wanted (down to 18), 5 offered (6 to 16).

\textbf{3.} Between 18 and 19 dalens: 5 baskets are wanted (30 to 20) and 5 offered (6 to 16). 5 baskets are traded.

\textbf{4.} The shoppers with values 30, 28, 25, 22 and 20; the farmers with costs 6, 8, 11, 13 and 16.

\textbf{5.} Pairing the highest values with the lowest costs: $24 + 20 + 14 + 9 + 4 = 71$; the sixth pair would give $18 - 19 = -1$. The largest \omterm{def:g10:surplus-and-efficiency:total-surplus}{total surplus} is
$125 - 54 = 71$ dalens, and the market's five trades reach it.

\textbf{6.} \omterm{def:g10:surplus-and-efficiency:consumer-surplus}{Consumer surplus} $125 - 5 \times 18.5 = 32.5$; \omterm{def:g10:surplus-and-efficiency:producer-surplus}{producer surplus} $5 \times 18.5 - 54 = 38.5$; total 71.

\textbf{7.} At 18: $125 - 90 = 35$ and $90 - 54 = 36$. At 19: $125 - 95 = 30$ and $95 - 54 = 41$. The total stays 71: the price only moves surplus from one side
to the other.

\textbf{8.} The basket is worth 18 to the shopper and costs 19 to produce: trading it would add $18 - 19 = -1$ dalen.

\textbf{9.} The first shopper gains $30 - 18.5 = 11.5$, the fifth $20 - 18.5 = 1.5$. No: efficiency is about which units are traded and by whom, and the price paid
moves surplus without changing the total.

\textbf{10.} 3 baskets, sold by the farmers with costs 6, 8 and 11.

\textbf{11.} Exactly three shoppers want a basket at prices above 22 and up to 25: the market settles between 22 and 25 dalens.

\textbf{12.} $(30 + 28 + 25) - (6 + 8 + 11) = 83 - 25 = 58$ dalens.

\textbf{13.} $71 - 58 = 13$ dalens, that is $13/71 \approx 18.3\,\%$ of the largest possible surplus.

\textbf{14.} The farmers with costs 13 and 16 lose their sales, and the shoppers with values 22 and 20 lose their baskets. The three remaining shoppers pay more
(22 to 25 instead of 18 to 19). The three farmers with stalls gain: at a price of 23.5, for example, their surplus is $3 \times 23.5 - 25 = 45.5$ against $3
\times 18.5 - 25 = 30.5$ without the cap.

\textbf{15.} The farmer with a cost of 26 sells only above 26, where only two shoppers buy. Two baskets are wanted at prices above 25 (up to 28) and two offered at
prices from 19 up to 26: the market settles between 25 and 26, with 2 baskets. \omterm{def:g10:surplus-and-efficiency:total-surplus}{Total surplus} $(30 + 28) - (11 + 19) = 28$ dalens.

\textbf{16.} $71 - 28 = 43$ dalens, $43/71 \approx 60.6\,\%$.

\textbf{17.} The stalls are held by dearer farmers: fewer baskets are worth trading, and those that are traded cost more to grow.

\textbf{18.} The farmer with a cost of 26 earns nothing from his stall, since he sells nothing at a price below 26. The farmer with a cost of 6 would earn about $25
- 6 = 19$ from it. Any payment between 0 and her gain leaves both better off.

\textbf{19.} With costs 6, 11 and 19, three baskets are traded (to the shoppers with values 30, 28 and 25), and \omterm{def:g10:surplus-and-efficiency:total-surplus}{total surplus} is $83 - 36 = 47$ dalens. The loss falls
to $71 - 47 = 24$: the trade recovered $47 - 28 = 19$ of the lottery's 43.

\textbf{20.} A cap of three stalls given to the cheapest farmers loses 13 dalens of surplus each Saturday, 18.3\,\% of the largest possible 71; given by lottery, the
same cap loses 43 dalens, 60.6\,\%.
\end{solution}
